\documentclass[11pt]{article}

\usepackage[margin=1in]{geometry}
\usepackage{amsmath,amssymb}
\usepackage{hyperref}
\usepackage{url}

\title{Free Energy as a Local Selection Rule in a\\ Sparse Finite-State Foraging Ecology}
\author{Alexander Kolpakov}
\date{}

\begin{document}
\maketitle

\begin{abstract}
We study whether free energy, in the loss-complexity form $F_\lambda = R + \lambda\,C$,
can serve as a \emph{local} selection rule over a population of deterministic agents,
and what such selection does to structure. The agents are sparse finite-state automata
(FSAs) that play a visible grid-foraging game; their genomes carry an explicitly
encoded transition relation over relational observations, and mutate by rule edits,
additions, and deletions. Selection minimizes $F_\lambda$, where $R$ is a foraging loss
and $C$ is a raw description length of the genome. Sweeping $\lambda$ traces a
loss-complexity landscape whose complexity variance behaves as a susceptibility-style
diagnostic. We report that fixed-task free-energy minimization reliably converges to
compressed policies, confirming the compression pole of the design, and we argue---but
do not yet demonstrate---that open-endedness requires solved structure to \emph{generate}
new validation pressure. We are explicit about which claims are established and which
are program.
\end{abstract}

\section{Introduction}
Open-ended evolution asks for a process that keeps producing qualitatively new structure
rather than settling into a fixed optimum. A recurring difficulty is that any single
fixed objective, pursued hard enough, is an attractor: the population converges and
stops innovating. Our working thesis is that free energy can be used not as a global
objective to be minimized once, but as a \emph{local} selection principle over agents
embedded in an expanding ecology---and that open-endedness then hinges on whether solved
structure produces new selection pressure. This paper builds the controllable instrument
required to test that thesis and reports what the closed-world version of it does.

We deliberately keep the substrate small and inspectable: deterministic automata, a
visible game, heritable and human-readable structure, replayable lineages, and an
explicit complexity term. The point is not a strong game bot; it is a phase-diagram
apparatus for free-energy selection.

\section{Background}
\paragraph{Description length and free energy.}
The complexity term $C$ is a raw description length in the sense of Kolmogorov
complexity~\cite{kolmogorov1965,live2008} and the minimum description length
principle~\cite{rissanen1978}, made concrete as the encoded size of a genome. Solomonoff
induction~\cite{solomonoff1964} supplies the prior intuition that simpler structures
should be preferred at equal fit. We treat the trade-off through the loss-complexity /
free-energy lens of Kolpakov~\cite{losscomplexity}: for a fitted model, $R$ (loss) and
$C$ (complexity) trace a landscape, and $\lambda$ selects a point on it. The functional
\[
  F_\lambda(a) \;=\; R(a) \;+\; \lambda\, C(a)
\]
is exactly a temperature-scaled free energy: at $\lambda \to 0$ selection is pure loss;
as $\lambda$ grows, parsimony dominates.

\paragraph{Open-endedness and self-improvement.}
Two threads inform the design. First, objective-free search: Lehman and
Stanley~\cite{lehman2011novelty} show that abandoning a fixed objective in favour of
novelty can outperform objective-driven search, motivating our shift from a global
objective to local selection plus frontier generation. Second, self-referential
improvement: Schmidhuber's G\"odel machine~\cite{schmidhuber2007godel} and
PowerPlay~\cite{schmidhuber2013powerplay} frame open-ended competence growth as
continually posing the simplest still-unsolved problem---the role we intend for an
archive of frontier environments.

\section{Method}
\paragraph{Agents.}
An agent is a sparse FSA. Its input is relational rather than raw pixels:
\[
  (\text{state},\ \text{previous\_move},\ \text{relative\_food\_azimuth})
  \;\longrightarrow\; (\text{move\_sequence},\ \text{next\_state}).
\]
A rule may emit a short \emph{sequence} of moves (a macro), not only a single move. If
no encoded rule matches the current input, the automaton \emph{halts} the episode: there
is no random fallback and no free default movement. This makes every behaviour traceable
to an explicit, paid-for rule.

\paragraph{Loss.}
$R(a)$ is a foraging loss: missed resources plus small per-step and collision costs. It
is bounded and deterministic given a map and a seed.

\paragraph{Complexity.}
$C(a)$ is a raw description length of the genome. We instrument four accounting modes to
test robustness of the conclusions to the definition of ``size'':
\emph{table} (default), the sum of complexities over the whole encoded sparse rule set,
so extra rules and long macros are paid for; \emph{active}, only rules exercised in
observed episodes; \emph{pruned}, rules reachable from the start state under any
observation; and \emph{mixed}, active complexity plus a dead-code tax on unused encoded
rules.

\paragraph{Selection and sweep.}
Selection minimizes $F_\lambda$ within an evolving population under mutation. The runner
sweeps $\lambda$ and records, per point, the loss, each complexity metric, and the
free energy, together with a \emph{complexity variance} across the population---a
susceptibility-style quantity we use to flag candidate transitions in the landscape.
Lineages and the best policy are exported for replay.

\section{Results}
On fixed maps (the closed-world condition) free-energy selection behaves as the
compression pole of the thesis predicts: loss falls quickly, then the population
compresses, shedding unused and redundant rules as $\lambda$ increases, and the
different complexity modes agree on the qualitative shape of the loss-complexity
frontier. The best policies are compact, deterministic, and replayable, and the
$\lambda$ sweep produces a monotone loss-complexity trade-off with an interior knee.
These results establish that (i) the substrate learns, (ii) $C$ is a real pressure and
not a tie-breaker, and (iii) the apparatus produces the diagnostics---frontier and
complexity variance---needed for the open-ended program.

\section{Limitations}
The results reported here are the \emph{closed-world} baseline: fixed maps, one static
loss. In that regime we observe exactly what a compression account predicts, which is
evidence for the mechanism but \emph{not} evidence of open-endedness. We do not yet
demonstrate solve--expand--solve cycles, an archive of frontier environments that
distinguishes lineages, or coevolutionary niche creation; those are designed but
unrealized. Complexity variance is presented as a diagnostic candidate, not a validated
order parameter. Nothing here should be read as a claim that free-energy selection
\emph{alone} yields open-endedness---our own thesis is that it does not, absent an
expanding ecology.

\section{Conclusion}
We built a small, inspectable instrument in which free energy $F_\lambda = R + \lambda C$
acts as a local selection rule over sparse deterministic foraging automata, and we
confirmed its compression behaviour and its diagnostics under a fixed task. The next
steps---frontier curricula near the competence boundary, archive validation, and
coevolution---are what would convert this from a compression demonstration into a test
of open-ended structural innovation.

\bibliographystyle{plain}
\bibliography{references}

\end{document}
